\section{Nearest-neighbor parents of injective PEPS}%
\label{sec:peps_injective_nearest_neighbor_parent}

\begin{definition}[Edge and isolated-site parent regions]%
    \label{def:peps_injective_nearest_neighbor_regions}
    \lean{TNLean.PEPS.edgePairRegion,TNLean.PEPS.edgeAndIsolatedRegion}
    \leanok
    For each edge $e=\{u,v\}$ of a finite graph, use the endpoint
    region $R_e=\{u,v\}$. The augmented family additionally contains
    the singleton $\{v\}$ for each isolated vertex $v$.
    All regional spaces are the images of the actual open-boundary
    contractions. Isolated-site conditions are needed to exclude
    unused physical directions at vertices absent from every edge;
    see \cite{gap:cpgsv21_injective_parent_reconstruction}.
\end{definition}

\begin{lemma}[Coverage by the nearest-neighbor regions]%
    \label{lem:peps_injective_nearest_neighbor_cover}
    \lean{TNLean.PEPS.edgePairRegion_covers_edges,
        TNLean.PEPS.mem_edgePairRegion_of_adj,
        TNLean.PEPS.edgePairRegion_covers_vertices,
        TNLean.PEPS.edgeAndIsolatedRegion_covers_vertices,
        TNLean.PEPS.edgeAndIsolatedRegion_covers_edges}
    \leanok
    \uses{def:peps_injective_nearest_neighbor_regions}
    Every edge is contained in its endpoint region. If every vertex
    has a neighbor, the endpoint regions also cover every vertex.
    On an arbitrary finite graph, the augmented family covers every
    vertex and every edge.
\end{lemma}

\begin{proof}\leanok
    A vertex with a neighbor belongs to the region of their edge.
    An isolated vertex belongs to its singleton region.
\end{proof}

\begin{theorem}[The nearest-neighbor parent space]%
    \label{thm:peps_injective_nearest_neighbor_parent_space}
    \lean{TNLean.PEPS.regionParentGroundSpace_edgePairs_eq_span,
        TNLean.PEPS.regionParentGroundSpace_edgeAndIsolated_eq_span,
        TNLean.PEPS.finrank_regionParentGroundSpace_edgePairs_eq_one,
        TNLean.PEPS.finrank_regionParentGroundSpace_edgeAndIsolated_eq_one}
    \leanok
    \uses{lem:peps_injective_nearest_neighbor_cover,
        thm:peps_injective_covering_parent_ground_space}
    Suppose every site tensor is injective.
    If every vertex has a neighbor, the common space for the
    edge-pair regions is exactly the span of the closed PEPS vector.
    For an arbitrary finite graph, the same conclusion holds for the
    augmented edge and isolated-site family.
    If every bond dimension is positive, each of these spaces is
    one-dimensional.
    These are the nearest-neighbor and isolated-site forms of the
    virtual-pair argument in
    \cite[Section~IV.C.1]{Cirac2021Matrix}.
    Without positive bond dimensions the span equality still holds,
    and the PEPS span may be zero.
\end{theorem}

\begin{proof}\leanok
    Apply the covering-region theorem to the two coverage statements.
    Positive bond dimensions make the virtual Bell product nonzero;
    injectivity preserves this property in the physical PEPS vector.
\end{proof}

\begin{corollary}[Exact-kernel nearest-neighbor Hamiltonians]%
    \label{cor:peps_injective_nearest_neighbor_parent_kernel}
    \lean{TNLean.PEPS.ker_regionParentHamiltonian_edgePairs_eq_span,
        TNLean.PEPS.ker_regionParentHamiltonian_edgeAndIsolated_eq_span,
        TNLean.PEPS.finrank_ker_regionParentHamiltonian_edgePairs_eq_one,
        TNLean.PEPS.finrank_ker_regionParentHamiltonian_edgeAndIsolated_eq_one}
    \leanok
    \uses{thm:peps_injective_nearest_neighbor_parent_space}
    Choose any positive semidefinite interaction on each region,
    with kernel exactly its actual regional PEPS space, and sum
    their extensions to the full physical graph.
    Under the corresponding hypotheses of the preceding theorem,
    the Hamiltonian kernel is exactly the closed PEPS span.
    Positive bond dimensions make this kernel one-dimensional.
    The edge-only conclusion requires that every vertex have a
    neighbor; the augmented conclusion includes the isolated-site
    singleton interactions explicitly.
\end{corollary}

\begin{proof}\leanok
    The kernel of a finite sum of positive interactions is the
    intersection of their kernels, which is precisely the common
    regional parent space.
\end{proof}
